\section{The One I Did Not Use}
\label{sec:ch09-the-one-i-did-not-use}

\index{provides@\code{"@provides}!what it is for|(}%
\code{@provides} is the third seam directive and this graph does not carry one.
That is a decision rather than an omission, and what it turns on is what the
directive costs to satisfy.

\index{provides@\code{"@provides}!is a routing hint, not a capability}%
What it does is narrow. A subgraph that returns an entity it does not own can
say that at this one field, it happens to have some of that entity's data
already, so the router need not make a second call to fetch it. Apollo's
reference is careful to describe it as
\enquote{a set of entity fields that a subgraph can resolve, but only at a
particular schema path}, and adds that using it
\enquote{is always an optional optimization}~\autocite{apollodirectives}. It
buys a round trip. It grants no new capability.

\index{provides@\code{"@provides}!has two preconditions and they are the point}%
The two preconditions are where the interest is. The providing subgraph must
declare the field and mark it \code{@external}, which is the same rule
\code{@requires} follows. And the field must be marked \code{@shareable} in at
least one other subgraph that defines it, because two services returning the
same field is exactly what \code{@shareable} exists to
permit~\autocite{apollodirectives}.

\index{provides@\code{"@provides}!needs the data to be in two places}%
Read those two together and the cost is plain. To provide \code{Speaker.name}
from the Sessions service, that service has to have speaker names, and the
Speakers service has to declare that it does not mind somebody else answering
for them. Which means keeping a copy of a column in a service that does not own
it, and writing the declaration that chapter~\ref{ch:composition} closes by
teaching you never to write.

\index{ownership!a saved round trip is not worth a duplicated column}%
That is the whole of my reason for not shipping one. This chapter spent its
first section moving speaker names out of the Sessions service because two
services claiming one field is the failure federation exists to prevent, and
\code{@provides} is a directive for the case where you have done it on purpose.
There are systems where the copy is already there for other reasons, and in
one of those the directive is free. Putting the copy back to save a fetch is
paying for a round trip with the thing chapter~\ref{ch:ownership} is about.

\subsection{What the Composer Actually Enforces}

\index{provides@\code{"@provides}!the composer does not check the second rule}%
Apollo's documentation says that if those preconditions are not met,
\enquote{a composition error occurs}~\autocite{apollodirectives}. That is
checkable, and I checked it, because a rule stated in prose and a rule enforced
by a tool are different things and this book has already met one pair that
disagreed.

\index{composition!an input written to fail the rule}%
The input is a hand-written document in the style
chapter~\ref{ch:composition} established: not a service, and the prose says so.
It provides \code{Speaker.name} at one field, marks it external locally, and is
composed against the real Speakers schema, where \code{name} carries no
\code{@shareable}:

\begin{graphqlsdl}
schema
  @link(
    url: "https://specs.apollo.dev/federation/v2.6"
    import: ["@key", "@external", "@provides", "FieldSet"]
  ) {
  query: Query
}

type Query {
  scheduleWithSpeakers: [Session!]!
}

type Session @key(fields: "id") {
  id: ID!
  speaker: Speaker @provides(fields: "name")
}

type Speaker @key(fields: "id") {
  id: ID!
  name: String! @external
}

scalar FieldSet
\end{graphqlsdl}

\index{composition!the second precondition is not enforced}%
By the documented rule that is a composition error. It composes:

\begin{minted}[fontsize=\footnotesize,breakanywhere]{text}
Router execution config successfully written to "router.json".
\end{minted}

\index{composition!vendors disagree about a rule, again}%
Exit zero, a config written, and nothing anywhere in the output mentioning
shareable. This is the second time in three chapters that this book has found
the tool and the specification's own documentation saying different things:
chapter~\ref{ch:composition} found Hot Chocolate emitting a \code{@cost}
directive whose argument Cosmo's composer types differently, and here is a rule
Apollo publishes and this composer does not apply.

\index{sources!which one to believe when they disagree}%
I am not going to tell you which of them is right, because I do not know and
the question has a real answer that neither of them states. What I can tell you
is which one will decide whether your deploy goes out, and it is the one you
run. Write the \code{@shareable} anyway if you use this directive: the
documented rule is the stricter of the two, a composer that starts enforcing it
would break a graph that composed yesterday, and there is no cost to being on
the safe side of a rule you were going to satisfy anyway.

\index{provides@\code{"@provides}!is weaker on an abstract type}%
One last thing, about where the weak ground is. Everything above concerns a
field whose type is an object. Put a \code{@provides} on a field returning an
interface or a union and you are in territory where implementations differ from
each other more than they do here, and I would test that case against your own
router before believing it rather than reasoning from what a plain object field
does.
\index{provides@\code{"@provides}!what it is for|)}%
